\section{Diagrama de flujo y pseudocódigo}

La lógica del Maestro está implementada como una máquina de estados no bloqueante. De forma resumida, la secuencia automática puede expresarse mediante el siguiente pseudocódigo:

\begin{enumerate}[leftmargin=*]
    \item Inicializar UART, LCD, I2C, VL53L0X y base de tiempo.
    \item Comprobar comunicación con los dos Slaves I2C.
    \item Colocar actuadores en estado seguro y mostrar el mensaje de bienvenida.
    \item Consultar periódicamente el sensor IR hasta obtener cinco detecciones consecutivas.
    \item Abrir la puerta de entrada y esperar 3 s.
    \item Inicializar el VL53L0X y arrancar el Stepper hacia adelante a 800 pasos/s.
    \item Cuando el VL53L0X se encuentre entre 265 y 270 mm, detener el Stepper y cerrar la puerta.
    \item Tomar una referencia con el HC-SR04, abrir el servo de agua y monitorear la distancia.
    \item Cuando la distancia ultrasónica aumente 20 mm, cerrar el agua y reanudar el Stepper.
    \item Cuando el VL53L0X se encuentre entre 125 y 130 mm, detener el Stepper y encender el motor DC con PWM 128.
    \item Mantener el cepillado durante 5 s, detener el motor DC y reanudar el Stepper.
    \item Cuando el VL53L0X mida 70 mm o menos, detener el proceso, asegurar actuadores y apagar el VL53L0X.
    \item Mover el Stepper en reversa durante 5 s y después detenerlo.
    \item Esperar a que el sensor IR indique que el vehículo fue retirado y regresar al estado de reposo.
    \item Durante los estados normales, comprobar periódicamente los enlaces I2C; ante un fallo, pasar al manejo de error y recuperación.
\end{enumerate}

El diagrama gráfico definitivo se incorporará posteriormente a partir de esta misma máquina de estados.
